\documentclass[10pt,a4paper]{amsart}
\usepackage{amsmath,amssymb,amsthm,graphicx}
% For commutative diagrams we use Xy-pic
\usepackage[all]{xy}

% We use 2cell for 2-commutative diagrams.
\xyoption{2cell}
% Theorem environments.
%
\theoremstyle{plain}
\newtheorem{theorem}[subsection]{Theorem}
\newtheorem{proposition}[subsection]{Proposition}
\newtheorem{lemma}[subsection]{Lemma}

\theoremstyle{definition}
\newtheorem{definition}[subsection]{Definition}
\newtheorem{example}[subsection]{Example}
\newtheorem{exercise}[subsection]{Exercise}
\newtheorem{situation}[subsection]{Situation}

\theoremstyle{remark}
\newtheorem{remark}[subsection]{Remark}
\newtheorem{remarks}[subsection]{Remarks}

\numberwithin{equation}{subsection}

% Macros
%
\def\lim{\mathop{\mathrm{lim}}\nolimits}
\def\colim{\mathop{\mathrm{colim}}\nolimits}
\def\Spec{\mathop{\mathrm{Spec}}}
\def\Hom{\mathop{\mathrm{Hom}}\nolimits}
\def\Ext{\mathop{\mathrm{Ext}}\nolimits}
\def\SheafHom{\mathop{\mathcal{H}\!\mathit{om}}\nolimits}
\def\SheafExt{\mathop{\mathcal{E}\!\mathit{xt}}\nolimits}
\def\Sch{\mathit{Sch}}
\def\Mor{\mathop{\mathrm{Mor}}\nolimits}
\def\Ob{\mathop{\mathrm{Ob}}\nolimits}
\def\Sh{\mathop{\mathit{Sh}}\nolimits}
\def\NL{\mathop{N\!L}\nolimits}
\def\CH{\mathop{\mathrm{CH}}\nolimits}
\def\proetale{{pro\text{-}\acute{e}tale}}
\def\etale{{\acute{e}tale}}
\def\QCoh{\mathit{QCoh}}
\def\Ker{\mathop{\mathrm{Ker}}}
\def\Im{\mathop{\mathrm{Im}}}
\def\Coker{\mathop{\mathrm{Coker}}}
\def\Coim{\mathop{\mathrm{Coim}}}

% Boxtimes
%
\DeclareMathSymbol{\boxtimes}{\mathbin}{AMSa}{"02}

%
% Macros for moduli stacks/spaces
%
\def\QCohstack{\mathcal{QC}\!\mathit{oh}}
\def\Cohstack{\mathcal{C}\!\mathit{oh}}
\def\Spacesstack{\mathcal{S}\!\mathit{paces}}
\def\Quotfunctor{\mathrm{Quot}}
\def\Hilbfunctor{\mathrm{Hilb}}
\def\Curvesstack{\mathcal{C}\!\mathit{urves}}
\def\Polarizedstack{\mathcal{P}\!\mathit{olarized}}
\def\Complexesstack{\mathcal{C}\!\mathit{omplexes}}
% \Pic is the operator that assigns to X its picard group, usage \Pic(X)
% \Picardstack_{X/B} denotes the Picard stack of X over B
% \Picardfunctor_{X/B} denotes the Picard functor of X over B
\def\Pic{\mathop{\mathrm{Pic}}\nolimits}
\def\Picardstack{\mathcal{P}\!\mathit{ic}}
\def\Picardfunctor{\mathrm{Pic}}
\def\Deformationcategory{\mathcal{D}\!\mathit{ef}}
\usepackage{cleveref}
\crefname{originalreference}{original reference}{original references}
\Crefname{originalreference}{Original reference}{Original references}
\makeatletter
\newlabel{item-base-change-prime-to-p}{{0A58}{source}{}{sourcepointer-1}{}}
\newlabel{item-base-change-prime-to-p@cref}{{[originalreference][0][]0A58}{[1][0][]source}{}{}{}}
\makeatother
\crefname{originalreference}{original reference}{original references}
\Crefname{originalreference}{Original reference}{Original references}
\makeatletter
\newlabel{item-base-change-proper}{{0A59}{source}{}{sourcepointer-2}{}}
\newlabel{item-base-change-proper@cref}{{[originalreference][0][]0A59}{[1][0][]source}{}{}{}}
\makeatother
\crefname{originalreference}{original reference}{original references}
\Crefname{originalreference}{Original reference}{Original references}
\makeatletter
\newlabel{item-finite-proper}{{0A57}{source}{}{sourcepointer-3}{}}
\newlabel{item-finite-proper@cref}{{[originalreference][0][]0A57}{[1][0][]source}{}{}{}}
\makeatother
\crefname{originalreference}{original reference}{original references}
\Crefname{originalreference}{Original reference}{Original references}
\makeatletter
\newlabel{item-surjective}{{06AP}{source}{}{sourcepointer-4}{}}
\newlabel{item-surjective@cref}{{[originalreference][0][]06AP}{[1][0][]source}{}{}{}}
\makeatother
\crefname{originalreference}{original reference}{original references}
\Crefname{originalreference}{Original reference}{Original references}
\makeatletter
\newlabel{item-vanishing}{{0A53}{source}{}{sourcepointer-5}{}}
\newlabel{item-vanishing@cref}{{[originalreference][0][]0A53}{[1][0][]source}{}{}{}}
\makeatother
\crefname{originalreference}{original reference}{original references}
\Crefname{originalreference}{Original reference}{Original references}
\makeatletter
\newlabel{lemma-affine-only-closed-points}{{09AY}{source}{}{sourcepointer-6}{}}
\newlabel{lemma-affine-only-closed-points@cref}{{[originalreference][0][]09AY}{[1][0][]source}{}{}{}}
\makeatother
\crefname{originalreference}{original reference}{original references}
\Crefname{originalreference}{Original reference}{Original references}
\makeatletter
\newlabel{lemma-cohomology-smooth-projective-curve}{{03RQ}{source}{}{sourcepointer-7}{}}
\newlabel{lemma-cohomology-smooth-projective-curve@cref}{{[originalreference][0][]03RQ}{[1][0][]source}{}{}{}}
\makeatother
\crefname{originalreference}{original reference}{original references}
\Crefname{originalreference}{Original reference}{Original references}
\makeatletter
\newlabel{lemma-compare-cohomology-point}{{03QU}{source}{}{sourcepointer-8}{}}
\newlabel{lemma-compare-cohomology-point@cref}{{[originalreference][0][]03QU}{[1][0][]source}{}{}{}}
\makeatother
\crefname{originalreference}{original reference}{original references}
\Crefname{originalreference}{Original reference}{Original references}
\makeatletter
\newlabel{lemma-finiteness-proper-variety-char-p-p}{{0A3P}{source}{}{sourcepointer-9}{}}
\newlabel{lemma-finiteness-proper-variety-char-p-p@cref}{{[originalreference][0][]0A3P}{[1][0][]source}{}{}{}}
\makeatother
\crefname{originalreference}{original reference}{original references}
\Crefname{originalreference}{Original reference}{Original references}
\makeatletter
\newlabel{lemma-vanishing-cohomology-mu-smooth-curve}{{03RR}{source}{}{sourcepointer-10}{}}
\newlabel{lemma-vanishing-cohomology-mu-smooth-curve@cref}{{[originalreference][0][]03RR}{[1][0][]source}{}{}{}}
\makeatother
\crefname{originalreference}{original reference}{original references}
\Crefname{originalreference}{Original reference}{Original references}
\makeatletter
\newlabel{lemma-vanishing-variety-char-p-p}{{0A3N}{source}{}{sourcepointer-11}{}}
\newlabel{lemma-vanishing-variety-char-p-p@cref}{{[originalreference][0][]0A3N}{[1][0][]source}{}{}{}}
\makeatother
\crefname{originalreference}{original reference}{original references}
\Crefname{originalreference}{Original reference}{Original references}
\makeatletter
\newlabel{situation-what-to-prove}{{0A52}{source}{}{sourcepointer-12}{}}
\newlabel{situation-what-to-prove@cref}{{[originalreference][0][]0A52}{[1][0][]source}{}{}{}}
\makeatother
\crefname{originallemma}{original lemma}{original lemmas}
\Crefname{originallemma}{Original lemma}{Original lemmas}
\makeatletter
\newlabel{varieties-lemma-curve-affine-projective}{{0A27}{source}{}{sourcepointer-13}{}}
\newlabel{varieties-lemma-curve-affine-projective@cref}{{[originallemma][0][]0A27}{[1][0][]source}{}{}{}}
\makeatother
\crefname{originallemma}{original lemma}{original lemmas}
\Crefname{originallemma}{Original lemma}{Original lemmas}
\makeatletter
\newlabel{varieties-lemma-proper-minus-point}{{0A24}{source}{}{sourcepointer-14}{}}
\newlabel{varieties-lemma-proper-minus-point@cref}{{[originallemma][0][]0A24}{[1][0][]source}{}{}{}}
\makeatother
\title{The vanishing and base change statements hold for the constant sheaf Z/ellZ on a smooth curve}
\author{The Stacks Project Authors}
\date{stacks-project:etale-cohomology.tex:local-snapshot-20261011; GNU FDL 1.2 or later}
\begin{document}
\maketitle
\noindent\textit{Editorial note.} All delivered mathematical text is the author's own, copied from the pinned Stacks Project chapter file etale-cohomology.tex; no statement or proof step was rewritten, reordered or removed.\par\medskip
\noindent\textit{Editorial note.} The counted proof is the author's complete argument: the proof is written out in full for both cases: the reduction to a curve, the projective and affine computations with the punctured-curve geometry in the prime-to-characteristic case, and the vanishing and finiteness lemmas in the characteristic case.\par\medskip
\noindent\textit{Editorial note.} The label lemma-constant-smooth-statements names the counted statement, and the proof environment immediately following it in the source is delivered verbatim as its proof.\par\medskip
\noindent\textit{Editorial note.} The counted statement is carried by the author's own lemma environment in the chapter source.\par\medskip
\noindent\textit{Editorial note.} The statement refers to the standing situation and to the enumerated statements by their labels; the fourteen cited inputs are kept as external pointers with their canonical Stacks tags.\par\medskip
\noindent\textit{Allowed standard prerequisites (editorial).}\begin{itemize}
\item Smooth schemes as disjoint unions of curves and points, and triviality on a point. Same chapter, the lemma comparing cohomology over a geometric point, and the lemma on the vanishing of the higher cohomology of an affine scheme whose points are closed. Supplies the reduction to the case of a smooth curve.
\item The projective and affine cases for coefficients prime to the characteristic. Same chapter, the lemma on the cohomology of a smooth projective curve and the lemma on the vanishing of the cohomology of mu on a smooth curve. Supplies the vanishing, finiteness and base change statements when ell is different from the characteristic.
\item The case of the characteristics and the geometry of punctured curves. Same chapter, the vanishing and finiteness lemmas for varieties over a field of characteristic p, and Varieties, the lemmas that a proper curve minus a point and an affine curve are related as stated. Supplies the characteristic-p case and the surjectivity statement through the affineness of a proper curve minus a point.
\end{itemize}\medskip
\ifdefined\hypertarget\hypertarget{sourcepointer-1}{}\fi
\noindent\textit{Source reference.} Original reference 0A58: sources/stacks-etale-cohomology/pinned-original-source.tex, source lines 15353--15353. Prerequisite pointer invoked at proof line(s) 15399; it names a standard fact of the same chapter that the counted argument uses, and every step of the counted proof built on it is retained verbatim in the delivered body.\par\medskip
\ifdefined\hypertarget\hypertarget{sourcepointer-2}{}\fi
\noindent\textit{Source reference.} Original reference 0A59: sources/stacks-etale-cohomology/pinned-original-source.tex, source lines 15359--15359. Prerequisite pointer invoked at proof line(s) 15411; it names a standard fact of the same chapter that the counted argument uses, and every step of the counted proof built on it is retained verbatim in the delivered body.\par\medskip
\ifdefined\hypertarget\hypertarget{sourcepointer-3}{}\fi
\noindent\textit{Source reference.} Original reference 0A57: sources/stacks-etale-cohomology/pinned-original-source.tex, source lines 15349--15349. Prerequisite pointer invoked at proof line(s) 15411; it names a standard fact of the same chapter that the counted argument uses, and every step of the counted proof built on it is retained verbatim in the delivered body.\par\medskip
\ifdefined\hypertarget\hypertarget{sourcepointer-4}{}\fi
\noindent\textit{Source reference.} Original reference 06AP: sources/stacks-etale-cohomology/pinned-original-source.tex, source lines 15365--15365. Prerequisite pointer invoked at proof line(s) 15381; it names a standard fact of the same chapter that the counted argument uses, and every step of the counted proof built on it is retained verbatim in the delivered body.\par\medskip
\ifdefined\hypertarget\hypertarget{sourcepointer-5}{}\fi
\noindent\textit{Source reference.} Original reference 0A53: sources/stacks-etale-cohomology/pinned-original-source.tex, source lines 15335--15335. Prerequisite pointer invoked at proof line(s) 15381; it names a standard fact of the same chapter that the counted argument uses, and every step of the counted proof built on it is retained verbatim in the delivered body.\par\medskip
\ifdefined\hypertarget\hypertarget{sourcepointer-6}{}\fi
\noindent\textit{Source reference.} Original reference 09AY: sources/stacks-etale-cohomology/pinned-original-source.tex, source lines 14280--14280. Prerequisite pointer invoked at proof line(s) 15390; it names a standard fact of the same chapter that the counted argument uses, and every step of the counted proof built on it is retained verbatim in the delivered body.\par\medskip
\ifdefined\hypertarget\hypertarget{sourcepointer-7}{}\fi
\noindent\textit{Source reference.} Original reference 03RQ: sources/stacks-etale-cohomology/pinned-original-source.tex, source lines 11135--11135. Prerequisite pointer invoked at proof line(s) 15396; it names a standard fact of the same chapter that the counted argument uses, and every step of the counted proof built on it is retained verbatim in the delivered body.\par\medskip
\ifdefined\hypertarget\hypertarget{sourcepointer-8}{}\fi
\noindent\textit{Source reference.} Original reference 03QU: sources/stacks-etale-cohomology/pinned-original-source.tex, source lines 9100--9100. Prerequisite pointer invoked at proof line(s) 15389; it names a standard fact of the same chapter that the counted argument uses, and every step of the counted proof built on it is retained verbatim in the delivered body.\par\medskip
\ifdefined\hypertarget\hypertarget{sourcepointer-9}{}\fi
\noindent\textit{Source reference.} Original reference 0A3P: sources/stacks-etale-cohomology/pinned-original-source.tex, source lines 9863--9863. Prerequisite pointer invoked at proof line(s) 15413; it names a standard fact of the same chapter that the counted argument uses, and every step of the counted proof built on it is retained verbatim in the delivered body.\par\medskip
\ifdefined\hypertarget\hypertarget{sourcepointer-10}{}\fi
\noindent\textit{Source reference.} Original reference 03RR: sources/stacks-etale-cohomology/pinned-original-source.tex, source lines 11246--11246. Prerequisite pointer invoked at proof line(s) 15398; it names a standard fact of the same chapter that the counted argument uses, and every step of the counted proof built on it is retained verbatim in the delivered body.\par\medskip
\ifdefined\hypertarget\hypertarget{sourcepointer-11}{}\fi
\noindent\textit{Source reference.} Original reference 0A3N: sources/stacks-etale-cohomology/pinned-original-source.tex, source lines 9838--9838. Prerequisite pointer invoked at proof line(s) 15410; it names a standard fact of the same chapter that the counted argument uses, and every step of the counted proof built on it is retained verbatim in the delivered body.\par\medskip
\ifdefined\hypertarget\hypertarget{sourcepointer-12}{}\fi
\noindent\textit{Source reference.} Original reference 0A52: sources/stacks-etale-cohomology/pinned-original-source.tex, source lines 15324--15324. Prerequisite pointer invoked at proof line(s) 15378; it names a standard fact of the same chapter that the counted argument uses, and every step of the counted proof built on it is retained verbatim in the delivered body.\par\medskip
\ifdefined\hypertarget\hypertarget{sourcepointer-13}{}\fi
\noindent\textit{Source reference.} Original lemma 0A27: sources/stacks-chow/tags.txt, source lines 12695--12695. Prerequisite pointer invoked at proof line(s) 15406; it names a standard fact of a sibling chapter that the counted argument uses, and every step of the counted proof built on it is retained verbatim in the delivered body.\par\medskip
\ifdefined\hypertarget\hypertarget{sourcepointer-14}{}\fi
\noindent\textit{Source reference.} Original lemma 0A24: sources/stacks-chow/tags.txt, source lines 12692--12692. Prerequisite pointer invoked at proof line(s) 15405; it names a standard fact of a sibling chapter that the counted argument uses, and every step of the counted proof built on it is retained verbatim in the delivered body.\par\medskip
\begin{lemma}
\label{lemma-constant-smooth-statements}
In Situation \ref{situation-what-to-prove}
assume $X$ is smooth and $\mathcal{F} = \underline{\mathbf{Z}/\ell\mathbf{Z}}$
for some prime number $\ell$. Then statements
(\ref{item-vanishing}) -- (\ref{item-surjective}) hold
for $\mathcal{F}$.
\end{lemma}

\begin{proof}
Since $X$ is smooth, we see that $X$ is a finite disjoint union of
smooth curves and points. The higher \'etale cohomology of any sheaf
on a point (the spectrum of $k$) is trivial, for example by
Lemma \ref{lemma-compare-cohomology-point}
or \ref{lemma-affine-only-closed-points}.
Hence we may assume $X$ is a smooth curve.

\medskip\noindent
Case I: $\ell$ different from the characteristic of $k$.
This case follows from
Lemma \ref{lemma-cohomology-smooth-projective-curve}
(projective case) and
Lemma \ref{lemma-vanishing-cohomology-mu-smooth-curve}
(affine case). Statement (\ref{item-base-change-prime-to-p})
on cohomology and extension of algebraically closed ground
field follows from the fact that the genus $g$ and the number
of ``punctures'' $r$ do not change when passing from $k$ to $k'$.
Statement (\ref{item-surjective}) follows as $H^2_\etale(U, \mathcal{F})$
is zero as soon as $U \not = X$, because then $U$ is affine
(Varieties, Lemmas \ref{varieties-lemma-proper-minus-point} and
\ref{varieties-lemma-curve-affine-projective}).

\medskip\noindent
Case II: $\ell$ is equal to the characteristic of $k$.
Vanishing by Lemma \ref{lemma-vanishing-variety-char-p-p}.
Statements (\ref{item-finite-proper}) and (\ref{item-base-change-proper})
follow from
Lemma \ref{lemma-finiteness-proper-variety-char-p-p}.
\end{proof}

\end{document}
